\begin{itemize}
	\item Se construyó una arquitectura que procesa predictores de reanálisis (ERA5) y satelitales (CHIRPS, MERRA2), incorporando reducción de dimensionalidad mediante PCA, transformación del objetivo y técnicas de regularización. La mejor combinación de predictores para la región resultó ser: la temperatura superficial del mar (SST), precipitación global (tp) e índices de la Oscilación Madden-Julian (RMM1, RMM2). Este resultado reduce la complejidad operativa y demuestra que variables de gran escala (SST y MJO) codifican suficiente información para pronosticar la precipitación de la región, sin necesidad de incluir otros campos complejos como geopotencial, viento o humedades relativa y específica.
	
	\item El modelo alcanzó un coeficiente de correlación de Pearson de \(r = 0.9166\), un coeficiente de Nash-Sutcliffe (\(\text{NSE}=0.7440\)) y un índice de Kling-Gupta (\(\text{KGE}=0.7503\)) sobre el conjunto de pronóstico retrospectivo. Estos valores indican una capacidad predictiva superior a la media climatológica y un equilibrio aceptable entre correlación, variabilidad y sesgo. No obstante, el modelo presenta sobreestimación (sesgo positivo de 47.22 mm, \(20.17\%\)), particularmente en los meses de inicio y pico de la estación lluviosa (abril, mayo, septiembre). El análisis mensual reveló mejor desempeño en junio-agosto (KGE \(>0.65\)) y dificultades en la estación seca (enero-febrero, diciembre) donde las métricas de eficiencia se vuelven negativas.
	
	\item  El modelo LSTM superó a la metodología NextGen en la mayoría de las métricas globales: menor RMSE (\(102.55\) vs \(114.15\) mm), mayor correlación de Pearson (\(0.9166\) vs \(0.8719\)) y mejor NSE (\(0.7440\) vs \(0.6931\)). NextGen presentó un sesgo porcentual menor (\(9.18\%\) frente a \(20.17\%\)) y un KGE ligeramente superior (\(0.8094\) vs \(0.7503\)). En términos espaciales, LSTM reprodujo con mayor fidelidad los gradientes finos de precipitación asociados a la orografía de la Bocacosta, mientras que NextGen produjo campos más suavizados. Ambos modelos enfrentan dificultades para pronosticar eventos extremos y durante la estación.
	
	\item El modelo LSTM desarrollado constituye una herramienta complementaria válida para el pronóstico de precipitación mensual en la región de Bocacosta, con la ventaja de requerir solo tres predictores de reanálisis (SST, tp, MJO) fácilmente disponibles. Su capacidad para capturar la variabilidad intraestacional y los patrones espaciales lo hace adecuado para sistemas de alerta temprana, planificación agrícola y gestión de recursos hídricos, así como un complemento al método de pronóstico NextGen ya utilizado en la institución.
\end{itemize}